\section{Security Properties}\label{sec:claims}

We state six security properties of the cA2A profile. Each is an informal invariant over the
runtime; formal proofs are left to future work. Each carries a status line for ca2a 0.3.1 that
names the layer at which it has been checked and what that check does not cover. This section
states the guarantees and mechanisms; results appear in Section~\ref{sec:evaluation}.

\begin{property}[Attenuation Soundness]
\label{prop:c1}
\emph{Status: validated in software.} For any delegation chain of signed credentials rooted at a
trusted issuer, the scope granted at each hop is a subset of the scope of its parent, and any
attempt to grant a capability the parent does not hold is rejected. This is the confused-deputy
defense \citep{confuseddeputy1988} within a chain: without it, a delegate could accept a narrow task and then
act with authority the delegator never conferred. The mechanism is a per-hop subset check in
\texttt{verify\_chain}: each credential carries a \texttt{scope} set, a \texttt{depth}, a
\texttt{parent\_id}, and a detached Ed25519 signature \citep{eddsa2017} over the canonical body.
A hop whose \texttt{scope} is not contained in its parent's \texttt{scope} raises
\texttt{ScopeEscalation}; a monotone-narrowing chain verifies. The semantics mirror
agent-manifest \citep{agentmanifest2026} and follow the capability-attenuation model of
macaroons \citep{macaroons2014}.
\end{property}

\begin{property}[In-Chain Replay and Splice Rejection]
\label{prop:c2}
\emph{Status: validated in software.} A credential presented twice within a chain, or lifted
from one chain and spliced into another, is rejected. This forecloses splicing a legitimately
issued grant into an unrelated delegation to inherit authority out of context. The mechanism
combines two checks in \texttt{verify\_chain}. First, each hop's \texttt{credential\_id} must be
unique within the chain; a repeat raises \texttt{CredentialReplay}. Second, structural
continuity binds each hop to the previous one: a hop's \texttt{parent\_id} must equal the
previous hop's \texttt{credential\_id} and its \texttt{issuer} must equal the previous hop's
\texttt{subject}. A credential spliced from a foreign chain breaks one of these links and raises
\texttt{BrokenDelegationLink}. The property does not cover presentation of a whole valid chain
by a party other than its delegate, since the verifier keeps no global history; on the live path
that case is refused by the holder proof (Section~\ref{sec:after}). Nor does it cover the delegate
itself presenting a valid chain with a request from a different task, since nothing binds a
chain to an execution (Section~\ref{sec:design}).
\end{property}

\begin{property}[Scope-Policy Intersection]
\label{prop:c3}
\emph{Status: validated in software, including the inbound path of the reference transport.}
The effective permission at an inbound peer call is the delegated scope intersected with the
receiving agent's local policy \citep{cedar2023}: a capability is exercised only if the upstream
chain granted it \emph{and} the local policy independently permits it. This makes each domain the
final authority over what runs on its own resources, so an over-broad delegated grant cannot force
action the receiver's policy forbids. The mechanism is the intersection of the verified chain
scope (Property~\ref{prop:c1}) with the local policy decision, evaluated in
\texttt{effective\_scope} and \texttt{enforce\_peer\_call} in \texttt{ca2a\_runtime.peer}; a
capability outside the intersection raises \texttt{ScopeNotPermitted}. Local policy is either a
Cedar policy set evaluated by \texttt{cedarpy} or a set-based \texttt{LocalPolicy}. The same
decision runs inside \texttt{handle\_peer\_request}, which the reference HTTP transport and the
\texttt{a2a-sdk} bridge call on each inbound request. The check does not cover third-party A2A
servers or a production deployment.
\end{property}

\begin{property}[Sealed-Payload Confidentiality]
\label{prop:c4}
\emph{Status: validated at the cryptographic layer.} A task payload sent across a trust boundary
decrypts only with the private key matching the channel key the caller appraised, so no party on
the connectivity path can read plaintext. The mechanism seals the payload to the peer's channel
public key with an HPKE-style X25519, HKDF-SHA256, and ChaCha20-Poly1305 construction
(\texttt{ca2a\_runtime.channel}); only the peer's private key opens it, an unrelated key fails,
and a tampered ciphertext fails closed under AEAD authentication. What software cannot establish
is that the private key is confined to the attested code, so that the operator hosting the peer
cannot read it. That depends on hardware appraisal being required on the call, on the
measurement covering the code that holds the key, and on platform protections such as SEV-SNP
ciphertext hiding. Section~\ref{sec:after} records that the measured hosts did not meet the last
two conditions. Until they are met, confidential payloads should not be sent across a trust
boundary in the belief that a compromised host cannot read them.
\end{property}

\begin{property}[Provenance Integrity]
\label{prop:c5}
\emph{Status: structural checks validated in software; authenticated lineage covered by the
project's tests.} The per-hop records for a workflow form a linked path that is bound to the
delegation authority. With signed TRACE records and trusted signing keys, tampering with any
record, reparenting a record, or attributing a record to a hop that did not sign it is detected
offline \citep{prov2013}. The mechanism adds a delegation-link block to each record:
\texttt{parent\_record\_hash} links a hop to its parent, and \texttt{credential\_id} names the
credential the hop acted under. Altering a record changes its hash and breaks the child's
\texttt{parent\_record\_hash} link, raising \texttt{ProvenanceLinkBroken}. A cross-check binds the
provenance path to the delegation chain of Property~\ref{prop:c1}: record $i$ must correspond to
credential $i$, and a signed record must carry the delegate key that credential names. Hash links
without signatures detect only partial edits: a writer who rewrites the whole path can recompute
every link. The property does not establish that declared work completed or that the submitted
path is complete.
\end{property}

\begin{property}[Cross-Operator Attestation]
\label{prop:c6}
\emph{Status: exercised with synthetic evidence; not validated on hardware across independent
operators.} Two agents in independent trust domains, holding independently generated keys, each
appraise the other's runtime attestation before a delegated task is accepted, and a peer that
loads a different binary than the one expected is detected. The mechanism follows the IETF RATS
architecture \citep{rats2021}: each peer produces a hardware-signed attestation report binding its
channel key and a nonce, and the counterpart verifies the report's certificate chain and
signature and compares its measurement with an expected value, so a binary swap changes the
measurement and is rejected. The SEV-SNP, TDX, and TPM verifiers that this requires ship in ca2a
0.3.1 (Section~\ref{subsec:attestation}). The recorded hardware runs are one-directional across
two clouds, or mutual within one operator's project; neither is mutual attestation between
independent operators.
\end{property}
